% from ega4-2-fr.tex
\oldpage[IV-2]{122}

\subsection{\texorpdfstring{Critères de cohérence pour les Modules
$\mathcal H^0_{X/Z}(\mathcal F)$}{Critères de cohérence pour les Modules H\string^0\string_\string{X/Z\string}(F)}}

\begin{proposition}[5.11.1]
\label{IV.2.s5.11.1-fr}
Soient $X$ un préschéma localement noethérien, $Z$ une partie de $X$ stable par
spécialisation, $\mathcal F$ un $\mathcal O_X$-Module cohérent. Désignons par
$(x_\alpha)$ la famille des points de
$\operatorname{Ass}(\mathcal F)\cap(X-Z)$ et, pour tout $\alpha$, soient
$Y_\alpha$ le sous-préschéma fermé réduit de $X$ ayant
$\overline{\{x_\alpha\}}$ pour espace sous-jacent,
$Z_\alpha=Z\cap\overline{\{x_\alpha\}}$. Alors les deux conditions suivantes
sont équivalentes:
\begin{enumerate}
\item[a)] $\mathcal H^0_{X/Z}(\mathcal F)$ est un $\mathcal O_X$-Module
cohérent.
\item[b)] Pour tout $\alpha$,
$\mathcal H^0_{Y_\alpha/Z_\alpha}(\mathcal O_{Y_\alpha})$ est un
$\mathcal O_{Y_\alpha}$-Module cohérent.
\end{enumerate}
Ces deux conditions entraînent la suivante:
\begin{enumerate}
\item[c)] Pour tout $\alpha$, on a
$\operatorname{codim}(Z_\alpha,Y_\alpha)\geq2$.
\end{enumerate}
De plus, les trois conditions a), b), c) sont équivalentes lorsqu'en outre une
des propriétés suivantes est vérifiée:
\begin{enumerate}
\item[(i)] Tout point de $X$ admet un voisinage ouvert isomorphe à un
sous-schéma d'un schéma régulier (auquel cas on dit aussi que $X$ est
localement immersible dans un schéma régulier).
\item[(ii)] Pour tout $\alpha$, $Y_\alpha$ est universellement caténaire
(5.6.2) et son normalisé (II, 6.3.8) $Y'_\alpha$ est fini sur $Y_\alpha$.
\end{enumerate}
\end{proposition}

Toutes les propriétés envisagées sont locales sur $X$, donc on peut se borner
au cas où $X=\operatorname{Spec}(A)$ est affine, $A$ étant un anneau
noethérien, et $\mathcal F=\widetilde M$, où $M$ est un $A$-module de type
fini. Alors, pour tout $\alpha$, si $h_\alpha$ est l'injection canonique
$Y_\alpha\to X$, $(h_\alpha)_*(\mathcal O_{Y_\alpha})=\mathcal G_\alpha$ est
le $\mathcal O_X$-Module correspondant au $A$-module quotient
$A/\mathfrak j_{x_\alpha}$, et, par définition de
$\operatorname{Ass}(\mathcal F)$, ce $A$-module est isomorphe à un
sous-$A$-module de $M$. Comme $\mathcal H^0_{X/Z}(\mathcal G_\alpha)$ est un
sous-$\mathcal O_X$-Module quasi-cohérent de
$\mathcal H^0_{X/Z}(\mathcal F)$ (5.9.2), l'hypothèse que
$\mathcal H^0_{X/Z}(\mathcal F)$ est cohérent entraîne qu'il en est de même de
$\mathcal H^0_{X/Z}(\mathcal G_\alpha)$. D'autre part, il résulte de la
définition (5.9.1.2) que $\mathcal H^0_{X/Z}(\mathcal G_\alpha)$ est
isomorphe à
$(h_\alpha)_*(\mathcal H^0_{Y_\alpha/Z_\alpha}(\mathcal O_{Y_\alpha}))$;
cela prouve que a) entraîne b). Pour voir que b) entraîne a), il suffit de
montrer qu'il y a une filtration finie $(\mathcal F_i)_{0\leq i\leq n}$ de
$\mathcal F$ formée de $\mathcal O_X$-Modules cohérents, avec
$\mathcal F_0=\mathcal F$, $\mathcal F_n=0$, et telle que
$\mathcal H^0_{X/Z}(\mathcal F_i/\mathcal F_{i+1})$ soit cohérent pour tout
$i$: cela résulte, par récurrence descendante sur $i$, des suites exactes
\[
0\longrightarrow\mathcal F_{i+1}\longrightarrow\mathcal F_i
\longrightarrow\mathcal F_i/\mathcal F_{i+1},
\]
du fait que $\mathcal H^0_{X/Z}$ est un foncteur exact à gauche (5.9.2), et
enfin de (0$_\mathrm I$, 5.3.3) et (I, 6.1.1). En vertu de (3.2.8), il suffit
donc de prouver que $\mathcal H^0_{X/Z}(\mathcal F)$ est cohérent lorsque
$\mathcal F$ est irredondant; autrement dit,
$\operatorname{Ass}(M)=\{\mathfrak p\}$ est réduit à un seul élément. Notons
maintenant le

\begin{lemma}[5.11.1.1]
\label{IV.2.s5.11.1.1-fr}
Soient $A$ un anneau noethérien, $M$ un $A$-module de type fini, tel que
$\operatorname{Ass}(M)=\{\mathfrak p\}$. Il existe une filtration finie
$(M_h)_{0\leq h\leq m}$ de $M$ telle que $M_0=M$, $M_m=0$ et que
$M_h/M_{h+1}$ soit isomorphe à un sous-module de $A/\mathfrak p$.
\end{lemma}

Notons d'abord que l'homomorphisme canonique $M\to M_\mathfrak p=N$ est
injectif (Bourbaki, \emph{Alg. comm.}, chap. IV, \S~1, n\textsuperscript{o}~2,
prop. 6). Posons $B=A_\mathfrak p$, $\mathfrak m=\mathfrak pA_\mathfrak p$,
idéal maximal de $B$; on a $\operatorname{Ass}(N)=\{\mathfrak m\}$
(\emph{loc. cit.}, prop. 5), et comme $N$ est un $B$-module de type fini,

